\section{可达、通信类与周期性}
\label{sec:06-Random-Processes/03-Markov-Chains/02}

拿到一个 Markov 链的转移矩阵，多数人的第一反应是解方程 \(\pi=\pi P\) 求平稳分布。但先别急——你得先搞清楚链能走到哪里。

设想一个转移图：状态 \(A\) 和 \(B\) 各自吸收（\(P_{AA}=P_{BB}=1\)），状态 \(C\) 以一半概率去 \(A\)、一半去 \(B\)。解 \(\pi=\pi P\) 会给你一族解——\(\delta_A,\delta_B\) 的任意凸组合都平稳。哪个才是``真正的''长期分布？取决于你从哪出发。

不先做结构分析就去解方程，你得到的答案可能不唯一，而你不理解它不唯一的原因。所以第一步不是代数，是图谱。

\subsection{可达与通信：画地图}

称状态 \(j\) 从 \(i\) 可达，如果存在某个步数 \(n\ge0\) 使

\[
(P^n)_{ij}>0,
\]

记作 \(i\to j\)。注意 \(n=0\) 对应单位矩阵——每个状态到自身总是可达的（零步即可）。

若 \(i\to j\) 且 \(j\to i\)，称二者通信，记作 \(i\leftrightarrow j\)。通信意味着你可以从 \(i\) 走到 \(j\)，也能从 \(j\) 走回来——虽然走的步数可能不同。

\subsection{通信把状态分成等价类}

通信关系是等价关系：
\begin{itemize}
\tightlist
\item
  自反性：\(i\leftrightarrow i\)（零步路径）；
\item
  对称性：定义本身对称；
\item
  传递性：若 \(i\to j\) 且 \(j\to k\)，拼接两条正概率路径即可 \(i\to k\)；反向同理。
\end{itemize}

因此通信将状态空间划分为互不重叠的通信类。每个类内部所有状态互相可达，类与类之间至多单向可达。

一个通信类若``进去了就出不来''——从类内任意状态出发不可能到达类外——称为闭类。单个吸收状态 \(i\)（满足 \(P_{ii}=1\)）是最简单的闭类。多个状态构成的闭类则像一个迷宫，有入口但无出口。

若整个状态空间只有一个通信类，链称为不可约。``不可约''不要求一步之内互相到达——只要求存在某条有限步数的正概率路径。

有限链的转移图提供直观工具：强连通分量（极大的互相可达子图）恰好对应通信类，没有指向外部边的分量就是闭类。你可以直接用图论算法分解转移矩阵。

\subsection{为什么先分类再求平稳分布}

如果链有多个闭类，长期行为取决于初始状态和随机路径：你最终会落入某个闭类，永远留在其中。每个闭类内部可能有自己的平稳分布，这些分布拼不出全空间唯一的平稳分布——任意凸组合也满足 \(\pi=\pi P\)。

两吸收态的例子就是极端情况：\(\delta_A\) 和 \(\delta_B\) 各自平稳，\(0.3\delta_A+0.7\delta_B\) 也平稳。方程告诉你的是一组解的线性结构，图告诉你的是这些解的来源——哪些状态会``赢''取决于初始条件和路上遇到的随机选择。

\subsection{周期：返回步数的最小公约数}

即使链不可约，行为也可能不是``越来越接近平稳''，而是永远振荡。

状态 \(i\) 的周期定义为

\[
d(i)
=\gcd\{n\ge1:(P^n)_{ii}>0\}.
\]

它是返回步数集合的最大公约数——你能回到 \(i\) 的步数都是它的倍数。若 \(d(i)=1\)，称状态非周期，可以在各种步数上返回。若 \(d(i)>1\)，返回被锁定在某种节拍上——比如只能偶数步回来。

在同一通信类中，所有状态周期相同。证明：取 \(m,n\) 使 \(i\to j\) 和 \(j\to i\) 有路径，从 \(i\) 到 \(i\) 的回路可以绕经 \(j\)，步数的 gcd 不变。因此可以谈一个不可约链的周期而不必指明哪个状态。

最经典的周期例子是两状态确定交替：

\[
P=
\begin{bmatrix}
0&1\\
1&0
\end{bmatrix}.
\]

从状态 1 出发，第 1 步必到 2，第 2 步必回 1，分布永远在 \((1,0)\) 和 \((0,1)\) 之间振荡。周期为 2。

这个链有平稳分布 \(\pi=(1/2,1/2)\)——但它``有''平稳分布不等于``收敛到''平稳分布。从状态 1 出发，分布永远不会接近 \((1/2,1/2)\)，而是永远在两个极端之间跳跃。这一区分在\hyperref[sec:06-Random-Processes/03-Markov-Chains/04]{``平稳分布、可逆性与遍历收敛''}一节深入讨论。

\subsection{自环是打破周期的最简单方法}

想让链非周期，最简单的办法是加自环。若存在状态 \(i\) 满足 \(P_{ii}>0\)，一步就可返回，周期必为 1，整个不可约链非周期。

在 MCMC 中常用``懒化''技巧：

\[
P_{\text{lazy}}=\frac12(I+P).
\]

每步以一半概率原地不动，一半概率按原矩阵转移。懒化后的链非周期（每个状态都加了自环），且保持原平稳分布：若 \(\pi P=\pi\)，则 \(\pi P_{\text{lazy}}=\frac12(\pi+\pi P)=\pi\)。代价是收敛变慢——你浪费了一半步数在踏步上。

但没有自环也可能非周期。如果一个状态既能两步返回又能三步返回，

\[
\gcd(2,3)=1,
\]

仍为非周期。关键在于返回步数集合的 gcd，而非单个步数是否等于 1。

\subsection{可约链中不能只报一个周期}

不同通信类可以有不同周期。一个吸收态周期为 1（你一直在那儿），另一个类可能周期为 3。暂留状态甚至可能根本没有返回路径——周期对它无定义。

不能对任意可约链笼统说``周期为二''而不指明说的是哪一部分。周期是针对通信类（或不可约子链）的概念，需要精确到类。

下一篇进一步追问：可达性告诉你存在一条正概率路径，但链真的会走回去吗？
